\documentclass[cjk]{loom}
\runningthread{Project 5: Pairs Trading}
% Cover follows the Project 5 Loom template; course design by Yuxiao Ma.
\definecolor{qpaper}{HTML}{F4F1E9}
\definecolor{qink}{HTML}{142D36}
\definecolor{qteal}{HTML}{247F82}
\definecolor{qgold}{HTML}{CF8951}
\newcommand{\coursecover}[2]{%
\hypersetup{pageanchor=false}
\begin{titlepage}\thispagestyle{empty}
\begin{tikzpicture}[remember picture,overlay]
\begin{scope}[shift={(current page.south west)},x=1cm,y=1cm]
\fill[qpaper] (0,0) rectangle (21,29.7);
\fill[qink] (0,27.5) rectangle (21,29.7);
\node[anchor=west,text=white,font=\sffamily\small] at (2,28.6) {QUANTITATIVE FINANCE / RESEARCH / SYSTEMATIC TRADING};
\node[anchor=north west,text=qink,align=left,font=\bfseries\fontsize{34}{46}\selectfont] at (1.95,25.1) {Quantitative Finance\\and Systematic Research};
\node[anchor=west,text=qteal,font=\fontsize{18}{24}\selectfont] at (2,20.5) {English Edition};
\draw[qteal,line width=1.6pt] (2,19.6)--(6.4,19.6);
\node[anchor=west,text=qink,font=\large] at (2,18.5) {From data and assumptions to verifiable research conclusions};
\node[anchor=north west,text=qink!70,align=left,font=\sffamily\small] at (2,15.7) {THEORY\\[6pt] MATHEMATICS\\[6pt] DATA\\[6pt] CODE\\[6pt] EVIDENCE};
\begin{scope}[shift={(13,11)}]
\foreach \r in {1,2,3,4}{\draw[qteal!30,line width=0.6pt] (0,0) circle (\r cm);}
\draw[qink!25] (-4.6,-3.6)--(-4.6,4.2); \draw[qink!25] (-4.6,-3.6)--(4.4,-3.6);
\draw[qteal,line width=1.4pt,domain=-2.6:3.4,samples=60,variable=\y] plot ({-2.9+0.45*sqrt(1+(\y-0.2)*(\y-0.2)/0.55)*2.2},{\y});
\draw[qgold,line width=1.2pt] (-4.6,-1.2)--(3.6,4.0);
\fill[qgold] (-1.36,0.86) circle (4pt);
\foreach \x/\y in {-1.0/-1.9,0.2/-0.6,1.1/0.4,-0.3/0.9,1.9/1.6,2.6/-1.0,0.8/2.4}{\fill[qink!75] (\x,\y) circle (2.2pt);}
\end{scope}
\draw[qink!20] (2,5.9)--(19,5.9);
\node[anchor=west,text=qink,font=\sffamily\fontsize{22}{28}\selectfont] at (2,4.7) {Yuxiao Ma};
\node[anchor=west,text=qink!65,font=\small] at (2,3.55) {#1};
\node[anchor=west,text=qink!55,font=\sffamily\small] at (2,2.65) {#2};
\node[anchor=east,text=qteal,font=\sffamily\bfseries\fontsize{46}{46}\selectfont] at (19,3.75) {Q};
\end{scope}
\end{tikzpicture}
\null\vfill
\end{titlepage}
\pagenumbering{roman}\hypersetup{pageanchor=true}\color{inkiron}}

\begin{document}
\color{inkiron}
\coursecover{Project 5 · Pairs Trading: Five Gates}{Fill-In Notes and Answers · BIL17B candidate · 16 October 2026}
\begin{strand}A pairs trade must pass five gates. Screening asks whether the spread may revert; hedging translates $\beta$ into dollars and shares; signals use only information then available; the ledger charges each leg correctly; evaluation reports losing as well as winning results. These notes follow the five gates and use roughly 70\% reading and 30\% fill-in work. Empirical tests on real prices are blocked pending data permission. All numbers come from hand calculations in the exercises or explicitly labelled constructed teaching panels; they demonstrate methods and are not market evidence.\end{strand}
\section*{The spine: five gates}
\footnotesize
\begin{tabularx}{\linewidth}{@{}p{1.5cm}p{2.8cm}L@{}}
\toprule
Gate & Question & \multicolumn{1}{l}{\color{indigo}Takeaway (fill in the blanks)}\\
\midrule
Screening & Will the spread revert? & \fillin[1.2cm] measures comovement; cointegration asks whether the \fillin[1.8cm] is stationary; test it against \fillin[1.8cm] critical values.\\
Hedging & How much to buy and sell? & Log $\beta$ is an \fillin[1.0cm]; the dollar ratio is $\beta$, and the absolute share ratio is \fillin[2.1cm].\\
Signal & When do we trade? & Estimate and freeze parameters on the \fillin[1.5cm]; a closing signal can first trade at the \fillin[2.0cm].\\
Ledger & Where does the cash go? & Spread and impact are in the \fillin[1.4cm]; separately deduct \fillin[1.0cm]. Charge interest and borrow fees using the balance and mark known at the interval \fillin[1.0cm]; \fillin[2.3cm].\\
Evaluation & What can we claim? & Compare with \fillin[1.4cm], report \fillin[2.3cm], and conclude under the preregistered rules.\\
\bottomrule
\end{tabularx}
\normalsize
\par\smallskip{\small Exercise mapping: P5-L01--L04 correspond to P5-01--P5-04; P5-L05 to P5-07; P5-L06 to P5-08; P5-L07 to P5-09; P5-L08 to P5-10; P5-L09 to P5-05; and P5-L10 to P5-06. Prerequisites: Chapters 16, 17, 29, 31, 35, and 36.}
\begin{yourturn}Cover the text and write one sentence for each gate: what goes in, and what comes out for the next gate.\workspace[5]\end{yourturn}
\input{sections/p05_s1_screen}
\input{sections/p05_s2_hedge}
\input{sections/p05_s3_signal}
\input{sections/p05_s4_ledger}
\input{sections/p05_s5_evaluate}
\input{sections/p05_answers}
\end{document}
